\documentclass[11pt]{article}
\usepackage{explorer}
\newcommand{\ExplorerLanguage}{FR}
\newcommand{\ExplorerNumber}{06}
\newcommand{\ExplorerKind}{Solutions détaillées}
\begin{document}
\ExplorerTitle{Lycée · graphes et logique}{Parcours sans lever le crayon}
\ExplorerSection{Exercice 1}
Un graphe a des degrés 2,3,3,4. Combien d’arêtes possède-t-il ?\par
La somme vaut 12, donc le graphe possède 6 arêtes.\par
\ExplorerSection{Exercice 2}
Prouver que le nombre de sommets de degré impair est pair.\par
La somme des degrés est paire. Les degrés pairs ne changent pas sa parité ; la somme des degrés impairs doit donc aussi être paire, ce qui impose un nombre pair de tels sommets.\par
\ExplorerSection{Exercice 3}
Peut-on tracer un carré et ses deux diagonales sans repasser sur une arête ? Les intersections intérieures ne sont pas des sommets.\par
Non : les quatre sommets ont degré 3. Quatre sommets impairs interdisent un parcours eulérien.\par
\ExplorerSection{Exercice 4}
Pourquoi deux cycles disjoints ne permettent-ils pas un parcours unique, alors que tous les degrés sont pairs ?\par
On ne peut passer d’une composante à l’autre par une arête du graphe. La condition de connexité est indispensable.\par
\ExplorerSection{Problème : raisonnement complet}
Un réseau a pour arêtes AB, BC, CD, DA, AC et CE ; toutes ont longueur 1 km. On veut partir de A, visiter chaque arête au moins une fois et revenir à A. Quels sommets sont de degré impair ? Quelle longueur minimale faut-il parcourir ? Donner un itinéraire atteignant cette longueur.\par
Les degrés de A,B,C,D,E sont 3,2,4,2,1 : A et E sont impairs. Un circuit fermé doit ajouter des passages reliant ces deux sommets ; leur plus courte liaison est A--C--E, de longueur 2 km. Le minimum est donc \(6 + 2 = 8\) km. L’itinéraire A--B--C--D--A--C--E--C--A l’atteint : AC et CE sont parcourues deux fois, les autres une fois.\par
\end{document}
